\subsection{PAT16: MRI-informed directional transport and static connectivity}\label{sec:pat16_methods}
A separate pilot reused preoperative sub-PAT16 derivatives from the BTC dataset and linked digital-twin workflow \citep{aerts2022,lundervold2026twin}. It used processed T1 anatomy, named CSF/gray-/white-matter probability maps, a lesion mask, a registered diffusion tensor, a cortical/subcortical atlas and resting-state BOLD with confounds. Existing fMRIPrep and MRtrix-derived inputs were read in place \citep{esteban2019,tournier2019}; preprocessing was not rerun. T1-Gd, T2 and FLAIR were not used. This is a hypothetical untreated 90-day evolution initialized from a preoperative lesion, not an estimate of tumor age, observed postoperative progression or patient growth rates.

With $u=c/K$, both variants solve $\partial_t u=\nabla\cdot(\mathbf D\nabla u)+\rho u(1-u)$. The isotropic reference uses $\mathbf D=d\mathbf I$, with separate assumed white-/gray-matter values. In eligible white matter, a unit principal dMRI axis $\bm v$ defines
\begin{equation}
 \mathbf D=D_\perp\mathbf I+(D_\parallel-D_\perp)\bm v\bm v^{\mathsf T},\qquad
 D_\perp=\frac{3d}{r+2},\quad D_\parallel=rD_\perp .\label{eq:pat16_tensor}
\end{equation}
Thus $\operatorname{tr}(\mathbf D)/3=d$ in both variants: measured water diffusion informs orientation, not the magnitude of cell migration. Positive tensor eigenvalues and $0.2<\mathrm{FA}\leq1$ define eligible white matter; other tissue remains isotropic. Tensor components are resampled and rigidly reoriented before eigendecomposition. The full 3 mm tissue domain uses conservative no-flux tensor transport, retaining mixed derivatives (Appendix~\ref{app:pat16}). Oxygen, necrosis, displacement and treatment are absent from this pilot.

We generated \PATScenarios{} paired Latin-hypercube scenarios (seed 16026), varying white-/gray-matter migration, proliferation, anisotropy, initial amplitude, lesion threshold and a small global directional perturbation. All distributions are explicit assumptions (Table~\ref{tab:pat16_priors}); they were not learned from PAT16 or transferred from synthetic recovery as patient estimates. Each pair shares its initial state and all rate draws. States were saved every 15 days through day 90. Outcomes are $\int u\,\dd V$ in mL of equivalent normalized burden, volume with $u\geq0.16$, and mean simulated density in atlas regions. None is a calibrated MRI tumor volume, absolute cell count or functional-risk measure. Empirical 5th/50th/95th percentiles describe prior/sensitivity propagation, not confidence or posterior credible intervals.

Static functional connectivity was recomputed from actual preprocessed BOLD after atlas resampling, censoring and joint nuisance regression. The 82 regions comprise 68 Desikan--Killiany cortical and 14 subcortical regions \citep{desikan2006}. Slice overlays project at most 35 strong correlations with $|r|\geq0.35$ whose ROI centers lie within a $\pm12$ mm slab. FC is not evolved or used to alter proliferation; it provides spatial context, not axonal paths or predicted functional loss. The executed notebook synchronizes time, plane, slice and density-quantile controls across isotropic, anisotropic, paired-difference and threshold-exceedance views.
